\section{Experimental Methodology}

\subsection{Experimental Parameters and Workloads}
To systematically map the performance frontier, RollupX is evaluated using a Python-based workload generator under three traffic profiles: steady-state (Poisson arrivals), bursty spikes (traffic cycling between low and high volume), and skewed distributions~\cite{silva2024zksyncdataset}. We perform sweeps across the following configuration variables:
\begin{itemize}
    \item \textbf{Batch Size (\(B_{\text{max}}\)):} The maximum number of transactions allowed in a single L2 batch before it is sealed.
    \item \textbf{Timeout (\(T_{\text{timeout}}\)):} The maximum time the sequencer waits before sealing and flushing a batch.
    \item \textbf{Sequencing Policies:} The rules used by the sequencer to order transactions in the mempool. We test standard First-Come-First-Served (FCFS), FeePriority (sorting by gas price), TimeBoost (sorting by bid price within discrete 5-second windows)~\cite{timeboost_docs}, and BlobPacking (prioritizing larger transaction sizes).
    \item \textbf{Data Availability (DA) Modes:} The storage path used for L2 transaction data on L1. We evaluate Calldata (writing data to L1 execution history), EIP-4844 Blobs (using cheaper, temporary blob storage)~\cite{eip4844}, and Offchain DA (publishing only state roots and ZK proofs while keeping transaction data off-chain).
    \item \textbf{Prover Backends:} The environment validating L2 state transitions. We compare a Mock Prover (simulating proving times using a configurable delay parameter) and the Real Prover (executing guest code in the RISC0 zkVM and wrapping STARK proofs into Groth16 SNARK proofs).
\end{itemize}

\subsection{Evaluation Metrics}
We collect telemetry logs across all lifecycle stages to compute five primary metrics:
\begin{itemize}
    \item \textbf{Committed TPS:} Realized transaction throughput settled on L1.
    \item \textbf{L2-to-L1 Latency:} Total time from transaction submission to L1 block inclusion.
    \item \textbf{Queue Wait Time (\(W_q\)):} Time spent in the sequencer mempool before sealing.
    \item \textbf{L1 Gas per Transaction:} Total L1 gas consumed divided by transaction count.
    \item \textbf{Proving Wall-Clock Time:} Prover proof-generation duration.
\end{itemize}
To evaluate latency fairness under different scheduling policies, we compute Jain's Fairness Index~\cite{motepalli2023sequencers}:
\begin{equation}
J(x_1, x_2, \dots, x_n) = \frac{\left(\sum_{i=1}^{n} x_i\right)^2}{n \sum_{i=1}^{n} x_i^2}
\end{equation}
where \(x_i\) represents the queue wait time of transaction \(i\).

\subsection{Analytical Dynamic Batching Model}
To formalize the trade-off between batch size, user wait times, and Layer-1 submission costs, we define a cost-penalty heuristic model~\cite{adler2021data, werner2024analyzing}. For a batch size \(N\), the penalty score \(U(N)\) is defined as:
\begin{equation}
U(N) = \frac{C_{\text{fixed}} + F_{\text{blob}} \left\lceil \frac{N b_{\text{tx}}}{B_{\text{blob}}} \right\rceil}{N} + \gamma \left( \frac{N}{2\lambda} + P(N) + T_{\text{L1}} \right)
\label{eq:dynamic_batching}
\end{equation}
where \(C_{\text{fixed}}\) is the fixed L1 verifier and bridge cost, \(b_{\text{tx}}\) is the transaction data size, \(B_{\text{blob}}\) and \(F_{\text{blob}}\) are blob size and cost, \(\lambda\) is the transaction arrival rate, \(P(N)\) is the prover time for batch size \(N\), \(T_{\text{L1}}\) is the L1 block mining delay, and \(\gamma\) is a weight factor balancing cost against latency.
